\chapter{Glossary}
\label{ch:glossary}

A quick reference to the acronyms and terms this book leans on. Where a term has
a chapter that treats it properly, the entry points there.

\begin{description}[style=nextline, leftmargin=0pt]

\item[ACATS] Ada Conformity Assessment Test Suite -- the official conformance
  tests for an Ada implementation. This project runs them on real silicon
  (Chapter~\ref{ch:acats-sweep}).

\item[ACK polling] Probing an I2C device's address until it answers. An EEPROM stops acknowledging while it programs a page, so the master polls instead of sleeping the worst case (Chapter~\ref{ch:eeprom}).

\item[ADU] \emph{Application Data Unit} -- a complete Modbus message as it travels: the \texttt{MBAP} header (TCP) or address+CRC (RTU) wrapping a \texttt{PDU} (Chapter~\ref{ch:modbus}).

\item[Alire] The Ada package manager (\texttt{alr}). It supplies the
  cross-toolchain and pins the runtime crate; no other SDK is needed.

\item[APB] \emph{Advanced Peripheral Bus} -- the lower-speed on-chip bus many ESP32-S3 peripherals are clocked from ($\sim$80\,MHz), divided down per peripheral.

\item[ARP] \emph{Address Resolution Protocol} -- resolves an IPv4 address to a \texttt{MAC} address on a local link; part of the W5500's hardwired stack.

\item[\texttt{bb-runtimes}] AdaCore's ``bareboard runtimes'' generator. The Ada
  runtime here is a board port grown from it (Chapter~\ref{ch:runtime-build}).

\item[BCD] \emph{Binary-Coded Decimal} -- each decimal digit in its own 4-bit nibble; the on-the-wire format of some RTC and sensor registers.

\item[\texttt{.bss}] The zero-initialised data segment; cleared by \texttt{start.S}
  at boot.

\item[CAN] \emph{Controller Area Network} -- a multi-drop automotive/industrial serial bus. The ESP32-S3's controller for it is the \texttt{TWAI} peripheral.

\item[CCOUNT / CCOMPARE] Xtensa cycle-counter and compare registers. The runtime
  uses them to drive the clock tick and timed delays.

\item[ceiling priority] The locking policy for protected objects: a caller is
  raised to the object's ceiling priority for the duration, bounding blocking and
  preventing priority inversion.

\item[closure] A subprogram bundled with the enclosing-scope variables it uses but
  does not declare -- the state it ``closes over.'' A nested subprogram that
  captures a local needs a stack \emph{trampoline}, which faults on this chip, so
  the HAL builds under \texttt{No\_Implicit\_Dynamic\_Code} and uses closure-free
  callbacks or tagged dispatch instead (\S\ref{driver:plugin}).

\item[context switch] Saving one task's CPU state and restoring another's. The
  Xtensa version (windows, the coprocessor, \texttt{THREADPTR}) is
  Chapter~\ref{ch:context-switch}.

\item[CRL] \emph{Certificate Revocation List} -- a CA-signed list of certificates revoked before expiry; an alternative to \texttt{OCSP} (Chapter~\ref{ch:tls}).

\item[CSPRNG] \emph{Cryptographically-Secure Pseudo-Random Number Generator} -- a \texttt{PRNG} whose output is unpredictable enough for keys and nonces.

\item[DAG] \emph{Directed Acyclic Graph} -- a graph with no cycles; e.g.\ the recursion-free call graph the stack analyser walks for a worst-case bound (Chapter~\ref{ch:static-mem}).

\item[DDR] \emph{Double Data Rate} -- transferring on both clock edges to double throughput, as in octal (\texttt{OPI}) PSRAM/flash timing.

\item[din tuning] Calibrating the read-data (\emph{data-in}) sampling phase --- the \texttt{din\_mode} knob --- that an octal PSRAM controller needs at 80\,MHz DDR. Its correct value is not knowable from source (it varies by chip, controller and temperature) and a wrong one hangs the CPU on a cache read with no recovery, so the bootloader measures it at boot with a bounded SPI1 sweep and takes the centre of the passing window (Chapter~\ref{ch:psram}).

\item[DRAM / IRAM] The data-bus and instruction-bus \emph{views} of the same
  internal SRAM. Code runs from the IRAM view; data lives in the DRAM view
  (\S\ref{ch:bootloader}).

\item[DROM / IROM] The data and instruction \emph{flash} windows, reached
  execute-/read-in-place through the cache MMU.

\item[ECC] \emph{Elliptic-Curve Cryptography} -- public-key cryptography over elliptic curves; smaller keys than RSA for equal strength (Chapter~\ref{ch:tls}).

\item[ECDH] \emph{Elliptic-Curve Diffie--Hellman} -- ECC key agreement; \texttt{ECDHE} is its ephemeral, forward-secret form (Chapter~\ref{ch:tls}).

\item[EEPROM] \emph{Electrically Erasable Programmable Read-Only Memory} -- small, byte-addressable non-volatile memory needing no erase before a write. The 24C family speaks I2C (Chapter~\ref{ch:eeprom}).

\item[ELF] \emph{Executable and Linkable Format} -- the object/executable file format GNAT emits, converted to a flash image for the board.

\item[FFS / FLS] \emph{Find-First-Set / Find-Last-Set} -- bit-scan operations the \texttt{TLSF} allocator uses to pick a size class in $O(1)$ (Chapter~\ref{ch:heap}).

\item[FIPS] \emph{Federal Information Processing Standards} -- US government security standards; some hardened servers are FIPS-configured and reject non-NIST curves (Chapter~\ref{ch:tls}).

\item[FSM] \emph{Finite-State Machine} -- a controller with a fixed set of states and transitions; several peripherals embed one in hardware.

\item[FTL] \emph{Flash Translation Layer} -- the layer that hides NOR flash's erase-before-write and limited endurance behind a plain read/write block interface, with wear levelling (Chapter~\ref{ch:storage}).

\item[GHASH] The keyed Galois-field hash that computes the authentication tag in AES-GCM (Galois/Counter Mode); the Ada half of the hybrid record cipher (Chapter~\ref{ch:tls}).

\item[GNARL] The GNU Ada Runtime Library -- the tasking half of the runtime
  (tasks, protected objects, the scheduler), archived as \texttt{libgnarl.a}.

\item[GNAT] The GNU Ada compiler. Its sequential runtime is \texttt{libgnat.a}.

\item[GNSS] \emph{Global Navigation Satellite System} -- the umbrella term for GPS, Galileo, BeiDou and the like.

\item[GPR / \texttt{gprbuild}] GNAT project files and their builder. Each crate
  (\texttt{esp32s3\_rts}, \texttt{esp32s3\_hal}) is a GPR project.

\item[HAL] Hardware Abstraction Layer -- the \texttt{ESP32S3.*} peripheral
  drivers (Chapters~\ref{ch:hal}--\ref{ch:hal3}).

\item[HKDF] \emph{HMAC-based Key Derivation Function} -- expands a shared secret into keying material; the heart of the TLS\,1.3 key schedule (Chapter~\ref{ch:tls}).

\item[ICMP] \emph{Internet Control Message Protocol} -- IP's diagnostic/error sublayer (ping lives here); part of the W5500's hardwired stack.

\item[IDF] \emph{(ESP-IDF) Espressif IoT Development Framework} -- the vendor's C SDK and HAL; this project deliberately does without it (Chapter~\ref{ch:encap}).

\item[IMU] \emph{Inertial Measurement Unit} -- an accelerometer plus gyroscope (sometimes magnetometer) motion sensor.

\item[IPI] Inter-Processor Interrupt -- the poke one core sends another to wake a
  task or rebalance, central to the SMP scheduler.

\item[ISRG] \emph{Internet Security Research Group} -- the nonprofit behind Let's Encrypt; its \emph{ISRG Root X1} CA is pinned by the HTTPS example (Chapter~\ref{ch:tls}).

\item[JEDEC] The semiconductor standards body; here, the JEDEC identifier every SPI-NOR flash returns (manufacturer + device id), read by the W25Q driver (Chapter~\ref{ch:storage}).

\item[Jorvik] A tasking profile (Ada RM~D.13) -- a less restrictive successor to
  Ravenscar that still bans the unbounded constructs. Both \texttt{light-tasking}
  and \texttt{embedded} build under it.

\item[LBA] \emph{Logical Block Addressing} -- a flat sector index into a block device; the unit the filesystem reads and writes (Chapter~\ref{ch:storage}).

\item[LRU] \emph{Least Recently Used} -- the eviction policy of the ext4 write-back block cache (Chapter~\ref{ch:storage}).

\item[LSB / MSB] \emph{Least-} / \emph{Most-Significant Bit} (or \emph{Byte}) -- which end of a value or transfer comes first; e.g.\ RGB565 pixels are sent MSB-first.

\item[MBAP] \emph{Modbus Application Protocol} header -- the seven-byte TCP framing (transaction, protocol, length, unit id) prefixing each \texttt{PDU} (Chapter~\ref{ch:modbus}).

\item[MCU] \emph{Microcontroller Unit} -- a CPU integrated with memory and peripherals; the ESP32-S3 itself.

\item[MMU] Memory Management Unit. Here it is the cache MMU that maps 64\,KB flash
  pages (and PSRAM) into the CPU's address windows.

\item[MPI] \emph{Multi-Precision Integer} -- an arbitrary-width integer; the workhorse of RSA and elliptic-curve big-number arithmetic (Chapter~\ref{ch:tls}).

\item[MSPI] The memory SPI controller that talks to flash and PSRAM; its timing
  is what makes PSRAM bring-up delicate (\S\ref{ch:bootloader}).

\item[NAT] \emph{Network Address Translation} -- rewriting addresses/ports at a gateway so many hosts share one public IP; the connect-out patterns here are NAT-friendly.

\item[NIST] \emph{National Institute of Standards and Technology} -- publishes the curves (P-256) and cipher/hash standards the TLS stack implements (Chapter~\ref{ch:tls}).

\item[NMEA] \emph{National Marine Electronics Association} -- its 0183 sentence format (e.g.\ \texttt{RMC}, \texttt{GGA}) is how GPS modules report a fix.

\item[OCSP] \emph{Online Certificate Status Protocol} -- a live query for whether a certificate has been revoked; an alternative to a \texttt{CRL} (Chapter~\ref{ch:tls}).

\item[OUI] \emph{Organizationally Unique Identifier} -- the top 24 bits of a \texttt{MAC} address, assigned per manufacturer.

\item[overloading] Several subprograms (or operators) sharing one name, told apart
  by their parameter and result \emph{profile}; the compiler resolves each call to
  the one whose profile matches. No keyword is involved -- distinct from
  \emph{overriding}. (E.g.\ \texttt{Send} overloaded on the CAN frame type,
  Chapter~\ref{ch:hal2}.)

\item[overriding] Replacing a tagged type's inherited primitive operation with a
  new body. The optional \texttt{overriding} keyword makes the compiler confirm the
  declaration really does replace an inherited operation -- catching a typo that
  would otherwise add a new one instead. The basis of dispatching, and the shape of
  the book's pluggable backends (override the handlers you support; the framework
  calls them by dispatch). Explained in Chapter~\ref{ch:ada-primer}; cf.\
  \emph{overloading}.

\item[PDM] Pulse-Density Modulation -- the 1-bit serial format some microphones
  and ADCs emit; the I2S driver converts it to/from PCM.

\item[PDU] \emph{Protocol Data Unit} -- the protocol-specific payload; for Modbus, a function code plus its data, carried inside an \texttt{ADU} (Chapter~\ref{ch:modbus}).

\item[PHY] The \emph{physical-layer} transceiver -- the front-end that turns link signalling into bits, e.g.\ the Ethernet PHY inside the W5500.

\item[PKCS] \emph{Public-Key Cryptography Standards} -- the RSA-era format family, e.g.\ PKCS\#1 signatures and PKCS\#8 keys (Chapter~\ref{ch:tls}).

\item[PLL] \emph{Phase-Locked Loop} -- the on-chip multiplier that synthesises the high CPU/peripheral clocks (e.g.\ 240\,MHz) from the crystal.

\item[PRNG] \emph{Pseudo-Random Number Generator} -- a deterministic generator of random-looking values; cf.\ \texttt{CSPRNG}.

\item[protected object] An Ada monitor: data plus operations with built-in mutual
  exclusion and condition synchronisation (\S\ref{ch:encap}).

\item[PSRAM] Pseudo-static external RAM, mapped at \texttt{16\#3D00\_0000\#} by the
  bootloader. Not DMA-capable.

\item[PSS] \emph{Probabilistic Signature Scheme} -- the modern randomised RSA signature padding (RSASSA-PSS), used for the TLS CertificateVerify (Chapter~\ref{ch:tls}).

\item[RAII] ``Resource Acquisition Is Initialization'' -- tying a resource's
  lifetime to an object's scope. In Ada it is a \texttt{controlled} type whose
  \texttt{Finalize} releases the resource (\S\ref{ch:encap}).

\item[Ravenscar] The original restrictive tasking profile for high-integrity
  real-time systems; Jorvik's stricter predecessor.

\item[rendezvous] Direct task-to-task synchronisation via \texttt{entry}/
  \texttt{accept}. Available only in the \texttt{full} profile.

\item[restriction profile] A named bundle of \texttt{pragma Restrictions} (e.g.
  Jorvik) that the compiler enforces, trading language features for analysability
  and footprint (Chapter~\ref{ch:profiles}).

\item[repeated START] A second I2C START issued without an intervening STOP, so the slave sees one command rather than two. It is how a master writes a register address and then reads that register (\texttt{ESP32S3.I2C.Write\_Read}, Section~\ref{driver:end-opcode}).

\item[RFC] \emph{Request For Comments} -- the IETF document series specifying the internet protocols (TCP, TLS, DHCP, FTP\,\dots).

\item[RTS] Run-Time System -- the compiled runtime an application links against;
  here, one of the three profile archives.

\item[SAN] \emph{Subject Alternative Name} -- the X.509 extension listing the hostnames a certificate is valid for; the leaf's SAN must cover the host (Chapter~\ref{ch:tls}).

\item[secondary stack] Where Ada returns unconstrained results (e.g. a
  \texttt{String} of run-time-determined length). Absent in \texttt{light-tasking}.

\item[SMP] Symmetric Multi-Processing -- both ESP32-S3 cores under one scheduler,
  with tasks pinned per core and woken across cores by IPI.

\item[SVD / \texttt{svd2ada}] A machine-readable register description (System View
  Description) and the AdaCore tool that turns it into typed Ada register packages
  (\S\ref{ch:runtime-build}).

\item[TDM] \emph{Time-Division Multiplexing} -- interleaving several channels on one line by time slot; the multi-slot I2S audio mode (Chapter~\ref{ch:hal2}).

\item[\texttt{THREADPTR}] An Xtensa per-thread pointer register, saved/restored on
  context switch so per-task thread-local data works.

\item[trampoline] A scrap of executable code the compiler writes onto the stack to
  carry a closure's captured state. The stack is non-executable here, so a
  trampoline faults; \texttt{No\_Implicit\_Dynamic\_Code} forbids generating one
  (\S\ref{driver:plugin}, Chapter~\ref{ch:debug}).

\item[VECBASE] The Xtensa register holding the base address of the exception/
  interrupt vector table; \texttt{start.S} points it at our vectors.

\item[watchdog (WDT)] A timer that resets the chip unless fed. The bootloader
  disables the boot watchdogs; the RTC super-watchdog is re-disarmed after a
  deep-sleep wake.

\item[window (register window)] The Xtensa calling convention rotates a window of
  physical registers per call instead of spilling to the stack; mismanaged
  windows are a recurring source of boot and context-switch bugs.

\item[XIP] Execute-In-Place -- running code directly from flash through the cache
  rather than copying it to RAM.

\item[ZCX] Zero-Cost Exceptions -- the exception model in \texttt{embedded} and
  \texttt{full}: no overhead unless an exception is actually raised.

\item[ZFP] ``Zero FootPrint'' -- a runtime (or coding style) with no tasking,
  exceptions or secondary stack; the 2nd-stage bootloader is written this way.

\end{description}
